\section*{Chapter \ref{ch:b3:many-electron-atoms} --- Many-Electron Atoms and Term Symbols}

\begin{solution}{exo:b3:many-electron-atoms:1}
\[
\Psi = \frac{1}{\sqrt6}\begin{vmatrix}1s\alpha(1) & 1s\beta(1) & 2s\alpha(1)\\
1s\alpha(2) & 1s\beta(2) & 2s\alpha(2)\\ 1s\alpha(3) & 1s\beta(3) & 2s\alpha(3)
\end{vmatrix}.
\]
$1s$ offers only two \omterm{def:b3:many-electron-atoms:spin-orbital}{spin-orbitals}, $1s\alpha$ and $1s\beta$; a third electron
would repeat one of them, making two columns equal: the determinant is zero.
\end{solution}

\begin{solution}{exo:b3:many-electron-atoms:2}
$\binom62 = 15$, $\binom63 = 20$, $\binom{10}{2} = 45$. $p^3$: ${}^4$S (4) +
${}^2$D (10) + ${}^2$P (6) = 20. $d^2$: ${}^3$F (21) + ${}^3$P (9) + ${}^1$G (9) +
${}^1$D (5) + ${}^1$S (1) = 45.
\end{solution}

\begin{solution}{exo:b3:many-electron-atoms:3}
N ($2p^3$, half full): \termsym{4}{S}{3/2}. O ($2p^4$, more than half):
\termsym{3}{P}{2}. F ($2p^5$): \termsym{2}{P}{3/2}. \ce{Ti^2+} ($3d^2$):
\termsym{3}{F}{2}. \ce{Fe^3+} ($3d^5$, half full, all spins parallel, $L = 0$):
\termsym{6}{S}{5/2}.
\end{solution}

\begin{solution}{exo:b3:many-electron-atoms:4}
${}^2$D: $J = \frac52$ (6 states) and $\frac32$ (4); $6 + 4 = 10 = 5 \times 2$.
${}^3$P: $J = 2$ (5), 1 (3), 0 (1); $5 + 3 + 1 = 9 = 3 \times 3$.
\end{solution}

\begin{solution}{exo:b3:many-electron-atoms:5}
Intervals $16.4$ ($J = 1 - 0$) and $27.0$ ($2 - 1$): ratio 1.65 instead of 2.
$A = 16.4/1 = \qty{16.4}{cm^{-1}}$ from the first, $27.0/2 = \qty{13.5}{cm^{-1}}$
from the second. A single $A$ does not fit: the levels are slightly mixed with
those of other terms (${}^1$D, ${}^1$S), and \omterm{def:b3:many-electron-atoms:russell-saunders}{Russell--Saunders coupling} is an
approximation.
\end{solution}

\begin{solution}{exo:b3:many-electron-atoms:6}
$\lambda = 10^7/\tilde\nu$: \qty{589.76}{nm} (D$_1$) and \qty{589.16}{nm}
(D$_2$). Separation $\qty{17.20}{cm^{-1}} = \qty{2.13}{meV}$. $E(\frac32) -
E(\frac12) = \frac32A$, so $A = \qty{11.47}{cm^{-1}}$.
\end{solution}

\begin{solution}{exo:b3:many-electron-atoms:7}
$3p \to 3s$: allowed ($\Delta l = -1$). $3d \to 3s$: forbidden ($\Delta l = -2$).
$3d \to 3p$: allowed. $4s \to 3s$: forbidden ($\Delta l = 0$, no change of
parity). $4s \to 3p$: allowed.
\end{solution}

\begin{solution}{exo:b3:many-electron-atoms:8}
Mean ${}^3$P: $(5 \times 169086.76 + 3 \times 169086.84 + 169087.83)/9 =
\qty{169086.91}{cm^{-1}}$. Gap to ${}^1$P: \qty{2047.98}{cm^{-1}} $=
\qty{0.254}{eV}$, so $K(1s,2p) = \qty{0.127}{eV}$, three times smaller than
$K(1s,2s)$. The \omterm{def:b3:many-electron-atoms:exchange}{exchange integral} measures the overlap of the two orbitals'
densities; the $2s$ orbital penetrates into the $1s$ region (it has density
near the nucleus), the $2p$ orbital vanishes at the nucleus and overlaps less.
\end{solution}

\begin{solution}{exo:b3:many-electron-atoms:9}
The largest $M_L = 4$ needs both electrons in $m_l = 2$, spins opposite: a
singlet, ${}^1$G (9 states). The next, $M_L = 3$, comes from $m_l = 2, 1$, with
$M_S = 1$ possible: ${}^3$F (21). Remaining $M_L = 2$ with $M_S = 0$ only:
${}^1$D (5). Remaining $M_L = 1$ with $M_S = 1$: ${}^3$P (9). One state left with
$M_L = M_S = 0$: ${}^1$S. Total 45.
\end{solution}

\begin{solution}{exo:b3:many-electron-atoms:10}
At $\mathbf r_1 = \mathbf r_2 = \mathbf r$ the triplet function is
$\frac{1}{\sqrt2}[1s(\mathbf r)2s(\mathbf r) - 2s(\mathbf r)1s(\mathbf r)] = 0$; the
singlet function is $\sqrt2\,1s(\mathbf r)2s(\mathbf r) \ne 0$. In the triplet the
electrons are never found at the same point and are on average farther
apart, so their repulsion is smaller: $E_+ - E_- = 2K > 0$.
\end{solution}

\begin{solution}{exo:b3:many-electron-atoms:11}
$d^8$ is more than half full: the coupling is inverted, $J = L + S = 4$
lowest. Intervals: $E(3) - E(4) = 1360.7$ and $E(2) - E(3) = 908.9$; ratio
1.50, against Landé's $4/3 = 1.33$. $|A| = 1360.7/4 = \qty{340}{cm^{-1}}$
($A < 0$). \omterm{def:b3:many-electron-atoms:spin-orbit}{Spin--orbit coupling} is much larger in this heavier ion than in
carbon (\qty{16}{cm^{-1}}).
\end{solution}

\begin{solution}{exo:b3:many-electron-atoms:12}
$\sum_J(2J+1)[J(J+1) - L(L+1) - S(S+1)] = 0$ (the $(2L+1)(2S+1)$ states can
be counted in either basis, and $\sum M_LM_S = 0$ over the uncoupled ones).
For ${}^3$P: $1(0 - 4) + 3(2 - 4) + 5(6 - 4) = -4 - 6 + 10 = 0$. Carbon's
weighted mean is $(0 + 3 \times 16.42 + 5 \times 43.41)/9 = \qty{29.6}{cm^{-1}}$:
the energy the term would have without \omterm{def:b3:many-electron-atoms:spin-orbit}{spin--orbit coupling}.
\end{solution}

\begin{solution}{pb:b3:many-electron-atoms:1}
\textbf{1.} $1s\alpha$, $1s\beta$, $2s\alpha$, $2s\beta$.
\textbf{2.} $|1s\alpha\,2s\alpha|$, $|1s\alpha\,2s\beta|$, $|1s\beta\,2s\alpha|$,
$|1s\beta\,2s\beta|$.
\textbf{3.} $|1s\alpha\,2s\alpha| = \frac{1}{\sqrt2}[1s(1)2s(2) - 2s(1)1s(2)]\,
\alpha(1)\alpha(2)$, by expanding the $2\times2$ determinant.
\textbf{4.} Their sum is $\frac{1}{\sqrt2}[1s(1)2s(2) - 2s(1)1s(2)]\cdot
\frac{1}{\sqrt2}[\alpha(1)\beta(2) + \beta(1)\alpha(2)]$ (times $\sqrt2$, then
normalised); their difference gives $\frac{1}{\sqrt2}[1s(1)2s(2) +
2s(1)1s(2)]\cdot\frac{1}{\sqrt2}[\alpha(1)\beta(2) - \beta(1)\alpha(2)]$.
\textbf{5.} Symmetric space with the antisymmetric spin function ($S = 0$);
antisymmetric space with the three symmetric spin functions ($S = 1$).
\textbf{6.} $2S + 1 = 1$ and 3: one and three states of the same energy.
\textbf{7.} $[\ce{Ar}]\,3d^2$; $\binom{10}{2} = 45$ determinants.
\textbf{8.} Maximum $S = 1$, maximum $L = 2 + 1 = 3$, less than half full:
\termsym{3}{F}{2}, the level at 0 in the data.
\textbf{9.} $21 + 9 + 9 + 5 + 1 = 45$.
\textbf{10.} Normal ($J = 2$ lowest). Intervals 184.9 and 235.5: ratio 1.27
against $4/3 = 1.33$; the rule holds within 5\,\%.
\textbf{11.} ${}^3$F: $(5 \times 0 + 7 \times 184.9 + 9 \times 420.4)/21 =
\qty{241.8}{cm^{-1}}$; ${}^3$P: $(10538.4 + 3 \times 10603.6 + 5 \times
10721.2)/9 = \qty{10661.7}{cm^{-1}}$.
\textbf{12.} $15B = \qty{10419.9}{cm^{-1}}$, $B = \qty{695}{cm^{-1}}$.
\textbf{13.} $\Delta S = 0$.
\textbf{14.} $\Delta S = 1$ is forbidden, and $1s2s \to 1s^2$ has $\Delta l = 0$
(no parity change; ${}^3$S$_1 \to {}^1$S$_0$ also has $\Delta L = 0$ between two S
states). The state is metastable: it lives far longer than ordinary excited
states.
\textbf{15.} $169086.91 - 159855.97 = \qty{9230.94}{cm^{-1}}$: \qty{1083.3}{nm}, in
the near infrared.
\textbf{16.} $171134.89 - 166277.44 = \qty{4857.45}{cm^{-1}}$: \qty{2058.7}{nm}.
\textbf{17.} $\qty{171134.89}{cm^{-1}}$: \qty{58.43}{nm}, in the vacuum
ultraviolet (absorbed by air).
\textbf{18.} Each family has its own lines and its own lowest state, and the
two never connect by light: they behave like two gases with two spectra.
\textbf{19.} $E_\pm = E_{1s} + E_{2s} + J \pm K$ (singlet $+$, triplet $-$).
\textbf{20.} The triplet: its spatial function vanishes when the electrons meet,
so they repel less.
\textbf{21.} \qty{6421.47}{cm^{-1}}, $6421.47/8065.54 = \qty{0.796}{eV}$.
\textbf{22.} Gap $\qty{2047.98}{cm^{-1}} = \qty{0.254}{eV}$, so $K(1s,2p) =
\qty{0.127}{eV}$.
\textbf{23.} $2s$ penetrates the $1s$ region and overlaps it more than $2p$
does.
\textbf{24.} Yes: in both configurations the triplet (larger $S$) lies lowest.
\textbf{25.} $K = \frac12 \times \qty{0.796}{eV}$: \textbf{the \omterm{def:b3:many-electron-atoms:exchange}{exchange
integral} $K(1s,2s)$ of helium is \qty{0.398}{eV}}, measured as half the
singlet--triplet gap.
\end{solution}
